\begin{frame}{확인 문제 (1)}
  \begin{itemize}
    \item \textbf{1.} 팀이 train 70\% / test 30\%만 쓰고
      하이퍼파라미터를 \textbf{test로} 고르겠다고 한다. 후보 10개,
      잡음 $\sigma=0.01$일 때 test 성능이 평균적으로 얼마나
      부풀어 오르는가?
      \vskip0.3em
      {\footnotesize \textbf{답}: Week 6.3 $\mathbb{E}[\max]$ 계산
      ($\sigma=1$일 때 $\approx 1.53$)에 $\sigma=0.01$을 곱하면
      $\approx 0.015$ --- test F1이 실제로 0.985인데 1.00으로
      보고될 수 있다. 질문 예: ``시드 5개로 val/test 순위의
      불일치율을 재볼 수 있나요?''}
    \item \textbf{2.} breast\_cancer(악성 62.7\%)에서 다수 클래스
      베이스라인의 val F1은 0.774였다. \textbf{정확도} 지표의
      베이스라인은? 왜 정확도가 F1보다 ``인정하기 쉬운''
      숫자인가?
      \vskip0.3em
      {\footnotesize \textbf{답}: 정확도 베이스라인 $\approx$ 62.7\%.
      정확도는 소수 클래스 몇 건을 잡았는지를 \textbf{희석}한다
      --- 50:50이 아니면 모델과 베이스라인의 차이를 과소/과대
      평가. F1은 희석 효과를 없애므로 적합.}
  \end{itemize}
\end{frame}
